% GENERATED FILE. Edit canonical lessons, then run npm run generate:books.
% canonical-source-hash: fnv1a64:b87fef35b449b173
% canonical-lessons: KA-C503-curuku, KA-C503-pratibhe, KA-C503-dadda, KA-C503-marevu, KA-C503-ekagrate

\chapter{Quick Minds and Slow}
\label{ch:ka-quick-minds-and-slow}
\hlchaptermodality{503}

\begin{chapteropening}
\textbf{By the end of this chapter:} \emph{I can say quick, talent, slow-witted, forgetfulness and concentration.}

Complete the last lesson of chapter 503: I can say quick, talent, slow-witted, forgetfulness and concentration.
\end{chapteropening}

\section[\texorpdfstring{\kn{ಚುರುಕು} (curuku)}{ಚುರುಕು (curuku)}]{\kn{ಚುರುಕು} (curuku) — quick; brisk}

\label{lesson:KA-C503-curuku}



\begin{quote}
Before the new one: say the Kannada for to take note, then the Kannada for to react.
\end{quote}

\subsection*{You'll want to know: \kn{ಚುರುಕು}}
\textbf{\kn{ಚುರುಕು}} — \emph{curuku} — \textquotedblleft{}quick; brisk\textquotedblright{}.

\begin{quote}
\kn{ಆ} \kn{ಹುಡುಗಿ} \kn{ಓದಿನಲ್ಲಿ} \kn{ತುಂಬಾ} \kn{ಚುರುಕು}.
\end{quote}

\begin{quote}
\emph{ā huḍugi ōdinalli tumbā curuku.}
\end{quote}

\begin{quote}
That girl is very quick at her studies.
\end{quote}

\kn{ಚುರುಕಾಗಿ}, briskly. Also used of a sharp sun: \kn{ಚುರುಕು} \kn{ಬಿಸಿಲು}.

\begin{cousinweb}
\textbf{\kn{ಚುರುಕು}} (\emph{curuku}) is native Dravidian: DEDR 2713 groups Kannada \emph{cuṟuku} with Tamil \emph{cuṟukku}, \textquotedblleft{}quickness, briskness, keenness\textquotedblright{}.
\end{cousinweb}

\begin{grammarlens}[title={What you've built}]
One more word for describing what you lived and felt.
\end{grammarlens}

\subsection*{Your turn}
\begin{itemize}
\raggedright
  \item \emph{Say it:} \emph{curuku}
  \item \emph{Say it:} \emph{curuku}, once more
  \item \emph{Say it:} \emph{ā huḍugi ōdinalli tumbā curuku.} — That girl is very quick at her studies.
  \item \emph{Recall:} say the Kannada for to take note, then the Kannada for to react, then say \emph{curuku} again
\end{itemize}

\begin{tcolorbox}[breakable,colback=teal!4,colframe=teal!35!black,title={Before you move on}]
What does \kn{ಚುರುಕು} mean? (\textquotedblleft{}Quick; brisk\textquotedblright{}.) Say it once more.
\end{tcolorbox}
\section[\texorpdfstring{\kn{ಪ್ರತಿಭೆ} (pratibhe)}{ಪ್ರತಿಭೆ (pratibhe)}]{\kn{ಪ್ರತಿಭೆ} (pratibhe) — talent; flair}

\label{lesson:KA-C503-pratibhe}



\begin{quote}
Before the new one: say the Kannada for to react, then the Kannada for quick.
\end{quote}

\subsection*{You'll want to know: \kn{ಪ್ರತಿಭೆ}}
\textbf{\kn{ಪ್ರತಿಭೆ}} — \emph{pratibhe} — \textquotedblleft{}talent; flair\textquotedblright{}.

\begin{quote}
\kn{ಆ} \kn{ಹುಡುಗನಿಗೆ} \kn{ಚಿತ್ರ} \kn{ಬಿಡಿಸುವ} \kn{ಪ್ರತಿಭೆ} \kn{ಇದೆ}.
\end{quote}

\begin{quote}
\emph{ā huḍuganige citra biḍisuva pratibhe ide.}
\end{quote}

\begin{quote}
That boy has a talent for painting.
\end{quote}

\kn{ಪ್ರತಿಭಾವಂತ}, a talented person.

\begin{cousinweb}
\textbf{\kn{ಪ್ರತಿಭೆ}} (\emph{pratibhe}) is borrowed from Sanskrit \emph{pratibhā}, \textquotedblleft{}light, splendour; intelligence, understanding\textquotedblright{} (Monier-Williams). The final -ā became -\kn{ೆ}.
\end{cousinweb}

\begin{grammarlens}[title={What you've built}]
One more word for describing what you lived and felt.
\end{grammarlens}

\subsection*{Your turn}
\begin{itemize}
\raggedright
  \item \emph{Say it:} \emph{pratibhe}
  \item \emph{Say it:} \emph{pratibhe}, once more
  \item \emph{Say it:} \emph{ā huḍuganige citra biḍisuva pratibhe ide.} — That boy has a talent for painting.
  \item \emph{Recall:} say the Kannada for to react, then the Kannada for quick, then say \emph{pratibhe} again
\end{itemize}

\begin{tcolorbox}[breakable,colback=teal!4,colframe=teal!35!black,title={Before you move on}]
What does \kn{ಪ್ರತಿಭೆ} mean? (\textquotedblleft{}Talent; flair\textquotedblright{}.) Say it once more.
\end{tcolorbox}
\section[\texorpdfstring{\kn{ದಡ್ಡ} (daḍḍa)}{ದಡ್ಡ (daḍḍa)}]{\kn{ದಡ್ಡ} (daḍḍa) — slow-witted; dull}

\label{lesson:KA-C503-dadda}



\begin{quote}
Before the new one: say the Kannada for quick, then the Kannada for talent.
\end{quote}

\subsection*{You'll want to know: \kn{ದಡ್ಡ}}
\textbf{\kn{ದಡ್ಡ}} — \emph{daḍḍa} — \textquotedblleft{}slow-witted; dull\textquotedblright{}.

\begin{quote}
\kn{ಗಣಿತದಲ್ಲಿ} \kn{ನಾನು} \kn{ಸ್ವಲ್ಪ} \kn{ದಡ್ಡ}.
\end{quote}

\begin{quote}
\emph{gaṇitadalli nānu svalpa daḍḍa.}
\end{quote}

\begin{quote}
I am a bit slow at mathematics.
\end{quote}

It can sting; \kn{ಜಾಣ} is the opposite.

\begin{cousinweb}
\textbf{\kn{ದಡ್ಡ}} (\emph{daḍḍa}) is native Dravidian: DEDR 2314 lists Kannada \emph{daḍḍa}, \textquotedblleft{}a dull, ignorant man\textquotedblright{}, beside Kodagu \emph{daḍḍë}, \textquotedblleft{}a dull man\textquotedblright{}.
\end{cousinweb}

\begin{grammarlens}[title={What you've built}]
One more word for describing what you lived and felt.
\end{grammarlens}

\subsection*{Your turn}
\begin{itemize}
\raggedright
  \item \emph{Say it:} \emph{daḍḍa}
  \item \emph{Say it:} \emph{daḍḍa}, once more
  \item \emph{Say it:} \emph{gaṇitadalli nānu svalpa daḍḍa.} — I am a bit slow at mathematics.
  \item \emph{Recall:} say the Kannada for quick, then the Kannada for talent, then say \emph{daḍḍa} again
\end{itemize}

\begin{tcolorbox}[breakable,colback=teal!4,colframe=teal!35!black,title={Before you move on}]
What does \kn{ದಡ್ಡ} mean? (\textquotedblleft{}Slow-witted; dull\textquotedblright{}.) Say it once more.
\end{tcolorbox}
\section[\texorpdfstring{\kn{ಮರೆವು} (marevu)}{ಮರೆವು (marevu)}]{\kn{ಮರೆವು} (marevu) — forgetfulness}

\label{lesson:KA-C503-marevu}



\begin{quote}
Before the new one: say the Kannada for talent, then the Kannada for slow-witted.
\end{quote}

\subsection*{You'll want to know: \kn{ಮರೆವು}}
\textbf{\kn{ಮರೆವು}} — \emph{marevu} — \textquotedblleft{}forgetfulness\textquotedblright{}.

\begin{quote}
\kn{ವಯಸ್ಸಾದ} \kn{ಮೇಲೆ} \kn{ಅಜ್ಜನಿಗೆ} \kn{ಮರೆವು} \kn{ಹೆಚ್ಚಾಗಿದೆ}.
\end{quote}

\begin{quote}
\emph{vayassāda mēle ajjanige marevu heccāgide.}
\end{quote}

\begin{quote}
With age, grandfather's forgetfulness has grown.
\end{quote}

The verb is \kn{ಮರೆ}, to forget; \kn{ಮರೆಯಬೇಡಿ}, please do not forget.

\begin{cousinweb}
\textbf{\kn{ಮರೆವು}} (\emph{marevu}) is native Dravidian: DEDR 4760 lists Kannada \emph{maṟe}, \textquotedblleft{}to forget, be forgotten\textquotedblright{}; -\emph{vu} makes the noun.
\end{cousinweb}

\begin{grammarlens}[title={What you've built}]
One more word for describing what you lived and felt.
\end{grammarlens}

\subsection*{Your turn}
\begin{itemize}
\raggedright
  \item \emph{Say it:} \emph{marevu}
  \item \emph{Say it:} \emph{marevu}, once more
  \item \emph{Say it:} \emph{vayassāda mēle ajjanige marevu heccāgide.} — With age, grandfather's forgetfulness has grown.
  \item \emph{Recall:} say the Kannada for talent, then the Kannada for slow-witted, then say \emph{marevu} again
\end{itemize}

\begin{tcolorbox}[breakable,colback=teal!4,colframe=teal!35!black,title={Before you move on}]
What does \kn{ಮರೆವು} mean? (\textquotedblleft{}Forgetfulness\textquotedblright{}.) Say it once more.
\end{tcolorbox}
\section[\texorpdfstring{\kn{ಏಕಾಗ್ರತೆ} (ēkāgrate)}{ಏಕಾಗ್ರತೆ (ēkāgrate)}]{\kn{ಏಕಾಗ್ರತೆ} (ēkāgrate) — concentration; focus}

\label{lesson:KA-C503-ekagrate}



\begin{quote}
Before the new one: say the Kannada for slow-witted, then the Kannada for forgetfulness.
\end{quote}

\subsection*{You'll want to know: \kn{ಏಕಾಗ್ರತೆ}}
\textbf{\kn{ಏಕಾಗ್ರತೆ}} — \emph{ēkāgrate} — \textquotedblleft{}concentration; focus\textquotedblright{}.

\begin{quote}
\kn{ಓದುವಾಗ} \kn{ನನಗೆ} \kn{ಏಕಾಗ್ರತೆ} \kn{ಬೇಕು}.
\end{quote}

\begin{quote}
\emph{ōduvāga nanage ēkāgrate bēku.}
\end{quote}

\begin{quote}
I need concentration while I study.
\end{quote}

\kn{ಏಕ}, one + \kn{ಅಗ್ರ}, point: a mind on one point.

\begin{cousinweb}
\textbf{\kn{ಏಕಾಗ್ರತೆ}} (\emph{ēkāgrate}) is borrowed from Sanskrit \emph{ekāgratā}, \textquotedblleft{}intentness in the pursuit of one object\textquotedblright{} (Monier-Williams).
\end{cousinweb}

\begin{grammarlens}[title={What you've built}]
One more word for describing what you lived and felt.
\end{grammarlens}

\subsection*{Your turn}
\begin{itemize}
\raggedright
  \item \emph{Say it:} \emph{ēkāgrate}
  \item \emph{Say it:} \emph{ēkāgrate}, once more
  \item \emph{Say it:} \emph{ōduvāga nanage ēkāgrate bēku.} — I need concentration while I study.
  \item \emph{Recall:} say the Kannada for slow-witted, then the Kannada for forgetfulness, then say \emph{ēkāgrate} again
\end{itemize}

\begin{tcolorbox}[breakable,colback=teal!4,colframe=teal!35!black,title={Before you move on}]
What does \kn{ಏಕಾಗ್ರತೆ} mean? (\textquotedblleft{}Concentration; focus\textquotedblright{}.) Say it once more.
\end{tcolorbox}
